\section{Background and related work}
\label{sec:background}

Tamper-evident logging and the normative mapping of AI regulation have matured separately. Logging research offers integrity that a third party can check, without reference to legal obligations; AI Act compliance research offers the mapping, without integrity that anyone but the producer can vouch for. This section reviews both lines and shows where they fail to meet.

\subsection{The AI Act's evidence-bearing obligations}

Regulation (EU) 2024/1689 (the AI Act)~\cite{ai_act} sets requirements for high-risk AI systems that can only be demonstrated through records. For the systems listed in Annex III, which include creditworthiness assessment, employment and access to essential services, the requirements of Chapter III apply from 2 December 2027. For high-risk systems under Art. 6(1) and Annex I they apply from 2 August 2028. Both dates were set by Regulation (EU) 2026/1744~\cite{omnibus_2026}, which amended Art. 113 AI Act.

That amendment changed little of what this article relies on. It extended simplified technical documentation to small mid-cap enterprises (Art. 11(1)), replaced the post-market monitoring template of Art. 72(3) with Commission guidance, and amended Art. 10(1), (5) and (6). It left untouched the provisions on bias examination and data quality used here, Art. 10(2)(f)--(g) and 10(3), and Arts. 12, 14, 15, 19, 73 and 78. Table~\ref{tab:obligations} lists the obligations that presuppose a record, what that record must show, and who reads it.

\begin{table*}[t]
\caption{AI Act obligations that presuppose a record, as amended by Regulation (EU) 2026/1744.}
\label{tab:obligations}
\footnotesize
\begin{tabularx}{\textwidth}{@{}>{\raggedright\arraybackslash}p{0.15\textwidth}LL@{}}
\toprule
Provision & Record that must exist & Who reads it \\
\midrule
Art. 10(2)(f)--(g), 10(3) & Examination of possible biases and the measures taken; relevance, representativeness and completeness of training, validation and test data & Provider; market surveillance authority on request (Art. 74(12)); notified body (Annex VII, 4.3) \\
Art. 11, Annex IV & Technical documentation kept up to date, including the data, metrics and validation procedures (Annex IV, 2(d), 2(g)) and the post-market monitoring plan (Annex IV, 9) & Market surveillance authority; notified body \\
Art. 12 & Automatic recording of events over the system's lifetime, sufficient to identify risk situations, support post-market monitoring and monitor operation & Provider; deployer (Art. 26(5)); authority on reasoned request (Art. 21(2)) \\
Art. 14 & Measures that let natural persons oversee the system, including deciding not to use, overriding or stopping it (Art. 14(4)) & Deployer; market surveillance authority \\
Art. 15 & Declared levels of accuracy and the metrics behind them (Art. 15(3)); robustness and cybersecurity measures & Deployer, through the instructions for use; market surveillance authority \\
Art. 19 & Retention by the provider of the automatically generated logs for at least six months & Provider; market surveillance authority \\
Art. 72 & Post-market monitoring that collects and analyses performance data, based on a plan & Provider; market surveillance authority \\
Art. 73 & Report of each serious incident within the deadlines of Art. 73(2)--(4) & Market surveillance authority \\
Art. 78 & Confidentiality of the information obtained, protecting trade secrets and source code & Binds the readers above \\
\bottomrule
\end{tabularx}
\end{table*}


The typical reader of this evidence is a market surveillance authority, not a notified body. For most Annex III systems the provider follows the internal control procedure of Annex VI (Art. 43(2)); notified bodies intervene only in the cases of Art. 43(1) and under the sectoral procedures for Annex I products. For credit institutions the authority is the financial supervisor (Art. 74(6)). Either way, the reader did not produce the evidence and must decide how far to trust whoever did. That is the problem this article addresses.

\subsection{Tamper-evident logging and ledger anchoring}

Making the alteration of a record detectable is an old problem with standard answers. Haber and Stornetta~\cite{haber1991} linked each timestamp to its predecessor through a hash, so that a back-dated insertion breaks the chain. Schneier and Kelsey~\cite{schneier1999} designed audit logs for an untrusted machine whose entries made before a compromise cannot be read or altered undetected. Crosby and Wallach~\cite{crosby2009} replaced the linear chain with a history tree built on Merkle trees~\cite{merkle1987}, giving logarithmic proofs that an entry is in the log and that one version of the log extends another.

Transparency logs put these structures into production. Certificate Transparency (RFC 9162)~\cite{rfc9162} publishes certificates in append-only Merkle logs whose signed tree heads let monitors detect a log that rewrites or forks its history. Sigstore applies the same design to software signing through its Rekor log~\cite{newman2022}. Trusted timestamping (RFC 3161)~\cite{rfc3161} adds a third party's attestation that data existed at a given time. None of these needs a blockchain: integrity follows from hashing, signatures and the publication of a commitment the log operator cannot retract.

Permissioned ledgers add governance over who may write, and in what order. Hyperledger Fabric~\cite{androulaki2018} separates execution, ordering and validation. Chaincode runs on endorsing peers, an ordering service sequences transactions into blocks, and each peer checks them against an endorsement policy before committing. Identities come from Membership Service Providers (MSPs), so every transaction names the organisation that submitted it. Fabric's ordering service uses Raft~\cite{ongaro2014}, which tolerates crash faults but not Byzantine ones; integrity against a faulty orderer rests on the endorsement policy and on each organisation keeping its own copy of the ledger.

Several works build audit trails on such ledgers. Regueiro et al.~\cite{regueiro2021} implement a blockchain audit trail that improves on the security and usability of conventional audit-log solutions. Lu et al.~\cite{lu2020} and Shi et al.~\cite{shi2022} build integrity-verification trails over Fabric, and Sychowiec and Zieliński~\cite{sychowiec2025} combine Fabric with Kafka for provenance across administrative domains. Ural and Yoshigoe~\cite{ural2023} survey blockchain-enhanced machine learning, in which recording training transactions on a ledger is presented as a route to a tamper-evident system. None of these works links its records to the legal obligations of an AI system.

The law has begun to recognise ledgers as evidence. Regulation (EU) 2024/1183~\cite{eidas_2024} amended eIDAS to make electronic ledgers a trust service: under Art. 45k(1), a ledger cannot be denied legal effect or admissibility as evidence solely because it is electronic or not qualified. Art. 45k(2) grants a presumption of unique and accurate sequential ordering and of integrity, but only to qualified electronic ledgers~\cite{alamillo2024}. Implementing Regulation (EU) 2025/2531~\cite{qel_2025} set their reference standards: ETSI EN 319 401 for every provider, and ISO 23257:2022 and ISO/TS 23635:2022 for providers that use distributed ledger technologies.

The point to retain is that tamper evidence does not require a blockchain. What a permissioned ledger adds to a signed log with an external anchor is a history that no single organisation can write or rewrite alone, and the identity of the organisation behind each entry. Whether that is worth its cost is an empirical question, and it motivates the arms of the ablation in Section~6.5 and the discussion in Section~7.2.

\subsection{Machine-readable compliance evidence}

NIST's Open Security Controls Assessment Language (OSCAL)~\cite{oscal} expresses compliance information as structured documents with published schemas. A catalogue lists controls, a profile selects and tailors them for a baseline, and an assessment-results document records observations, findings against control objectives and an assessment log. OSCAL was built for security frameworks such as NIST SP 800-53; NIST publishes no catalogue for the AI Act.

Work aimed at the AI Act has produced the mapping but not the integrity. Ugarte et al.~\cite{ugarte2026} use OSCAL for machine-readable compliance and compute SHA-256 digests of the data splits and model weights, but nothing signs or anchors the OSCAL document itself. Buscemi et al.~\cite{buscemi2026} map AI Act requirements onto controls and tests on a real automotive system, and prescribe tamper-evident logs among those controls; the methodology leaves no signed record of a verification when it occurs. TechOps~\cite{lucaj2025} and Compliance Cards~\cite{marino2024} provide structured documentation aligned with Annex IV and computational compliance artefacts, without verifiable integrity.

Two works come closer. Palumbo et al.~\cite{palumbo2025} propose a governance-oriented observability framework for the AI Act and name, as future work, anchoring its records in a tamper-evident ledger; this article takes that step. Ramos and Ellul~\cite{ramos2024} propose blockchain for AI Act compliance from a cybersecurity angle, at a conceptual level and without implementation.

The antecedents are older. FactSheets~\cite{arnold2019} anticipated Annex IV-style documentation as a supplier's declaration of conformity for AI services, and Hacker et al.~\cite{hacker2022} framed AI compliance as a field that bridges data science and law. In industry, policy-as-code engines such as Open Policy Agent~\cite{opa} evaluate declarative rules against structured input at run time. The engine described in Section~5.2 follows that pattern, and records each decision as signed evidence.

\subsection{Data-protection constraints}

A hash is not outside the GDPR by construction. Personal data is any information relating to an identifiable person (Art. 4(1) GDPR~\cite{gdpr}), and Recital 26 requires account of all the means reasonably likely to be used to identify them. Finck~\cite{finck2018} argues accordingly that hashing does not amount to anonymisation: a hash that can be correlated with external identifiers remains personal data.

The European Data Protection Board takes the same position. Its Guidelines 02/2025 on blockchain (version 2.0, adopted on 7 July 2026)~\cite{edpb_2026} treat the hash of personal data as personal data and state that the GDPR still applies when only the hash is stored on chain (para. 52). They advise keeping personal data off chain wherever possible and prefer keyed or salted hashes to plain ones. Hash-on-chain is therefore a minimisation measure, not an exemption.

Redactable blockchains~\cite{ateniese2017} and Fabric's private data collections, whose private payload can be purged while its hash stays on the channel, allow limited erasure. This article uses neither: it keeps personal data out of the evidence record altogether (Section~5.4).

\subsection{Positioning}

Table~\ref{tab:positioning} scores the reviewed works on six criteria, with the same rule for every row, including this article's. A work scores ``partial'' when it addresses a criterion in design but not in an artefact, or for part of its scope.

\begin{table*}[t]
\caption{Positioning against related work. C1: explicit mapping to AI Act obligations. C2: integration with a running ML pipeline. C3: third-party-verifiable integrity (signature or anchoring). C4: statement of evidentiary limits. C5: controlled disclosure. C6: measured cost. A work scores ``partial'' when it addresses a criterion in design but not in an artefact, or for part of its scope.}
\label{tab:positioning}
\footnotesize
\begin{tabularx}{\textwidth}{@{}Lcccccc@{}}
\toprule
Work & C1 & C2 & C3 & C4 & C5 & C6 \\
\midrule
Schneier and Kelsey~\cite{schneier1999}; Crosby and Wallach~\cite{crosby2009} & no & no & yes & partial & partial & partial \\
Certificate Transparency~\cite{rfc9162}; Sigstore~\cite{newman2022} & no & partial & yes & partial & no & CHECK \\
Regueiro et al.~\cite{regueiro2021} & no & no & yes & no & no & CHECK \\
Shi et al.~\cite{shi2022} & no & no & yes & no & partial & CHECK \\
Ural and Yoshigoe~\cite{ural2023} & no & partial & partial & no & no & no \\
Ramos and Ellul~\cite{ramos2024} & yes & no & partial & no & no & no \\
Marino et al.~\cite{marino2024} & yes & partial & no & no & no & no \\
Palumbo et al.~\cite{palumbo2025} & yes & yes & no & no & no & no \\
Lucaj et al.~\cite{lucaj2025} & yes & partial & no & no & no & no \\
Ugarte et al.~\cite{ugarte2026} & yes & partial & partial & no & no & no \\
Buscemi et al.~\cite{buscemi2026} & yes & partial & no & no & no & no \\
This article & partial & yes & yes & yes & partial & yes \\
\bottomrule
\end{tabularx}
\end{table*}


This article scores ``partial'' on C1 because it covers seven provisions, not the whole of Chapter III, and on C5 because its redaction works by field and does not prevent inference by aggregation (Section~5.5).

Read by column, the table shows two families with complementary gaps. Logging and ledger work offers third-party-verifiable integrity, and some of it measures its cost, but none maps its records to AI Act obligations. AI Act work maps obligations to controls and documents, yet leaves their integrity to trust in the producer; and among the reviewed works only the classic logging literature states, through its threat models, what its records do not establish. This article joins the two lines and makes that statement explicit for each obligation (Section~4).
